\section{One ML/AI analytics use case}
AI1 detects a person crossing a line towards the car park in the vehicle lane. C08 is at SITE $(63,20,3)$, on the inside of the lane's northwest corner. Settings are azimuth 292$^\circ$, tilt 25$^\circ$, HFOV 72$^\circ$, VFOV about 44.46$^\circ$, and 1920 $\times$ 1080 without roll or digital cropping. WDR and distortion correction are enabled; the projection and border reserve defined in the CCTV section apply.\cite{S-CAM,S-AOA}

The event line runs from $(63.3,15.5)$ to $(66.7,15.5)$, south to north. The feet-detection region is $x=63.3$--66.7 and $y=12.5$--16.5. The outer 0.3 m strips by the mesh are outside the AI claim and remain under patrol observation. The whole-lane context area $x=63$--67, $y=10$--20 is not the analysis region. The separate pedestrian route lies outside it. An event indicates movement at a specified location, without assessing intention, authorisation, or identity.

Object Analytics places a person's detection point at the feet. Checks therefore include $z=0$ and the whole-person envelope up to 2 m high, extending 0.3 m each way in plan. They cover both ends and the entire line, plus corners and a 0.25 m regional grid. The envelope and camera rays remain inside the convex lane $x=63$--67, $y=10$--20: they do not cross side mesh or closed G01 at $y=10$; the envelope's lower plan edge at $y=12.2$ also clears open leaves ending at $y=12$. Vehicles or other people may still temporarily obscure the target.\cite{S-AOA}

For a minimum person height of 1.4 m, the smallest calculated image height is \AIHeightPx{} px against the recommended 8\% of image height, or 86.4 px. An object walking at up to 1.5 m/s has 3 m of visible path before the line, giving at least 2 s in the scene. This is a scenario design limit; equal performance for running or shorter people is not promised. The export also records the feet polygon in image coordinates for rule configuration.\cite{S-AOA}

AXIS Object Analytics classifies people and tracks line crossing on C08's DLPU (ARTPEC-8). Video input is 25 fps and analysis runs continuously in the camera. The manufacturer publishes no guaranteed inference rate for this scene, so 25 fps is not equated with 25 analyses/s. Actual event delay and processing slowdowns are measured at acceptance. No second AI server or eight-stream analytics installation is provided.\cite{S-CAM,S-AOA}

The native Axis event reaches the VMS rule on SOC-0, displays C08, and bookmarks a clip 10 s before / 20 s after the event. The operator acknowledges it and calls patrol after assessing the image. Repeated events for the same object are merged over 10 s. Analytics does not control doors, the gate, or FAS. Missing video does not automatically cancel another system's alarm.

Acceptance requires at least 100 positive scenarios, 100 negative scenarios, and 40 camera-hours across different days and lighting conditions. Tuning and test sets are separated by day. An independent observer labels crossing and time; unscorable cases are described separately, not silently removed from the report.
\[
\mathrm{precision}=\frac{TP}{TP+FP},\qquad
\mathrm{recall}=\frac{TP}{TP+FN}.
\]
Project targets are precision and recall of at least 0.90, at most 0.1 false events per camera-hour, and p95 delay of up to 2 s from line crossing to SOC alarm display. With 60 s handling per false event, this corresponds to 0.1 min of operator work per hour. True events and total shift workload are reported separately. The report includes TP/FP/FN denominators, uncertainty intervals, weather, occlusion, and lost events during network outages. None of these thresholds is a completed experiment result.
